\section{What cut elimination establishes}
\label{sec:cut-elimination-02}

Cut elimination establishes admissibility of cut and yields analytic
derivations.  Depending on the calculus, it supports consistency,
interpolation, conservativity, and proof-search results.  None of those
consequences follows merely from the phrase ``cut-free'': each requires a
separate argument using the exact rule set.

Cut elimination does not imply that removing cut makes proof search efficient.
Elimination can cause a non-elementary increase in derivation size, and a
cut-free search space may still be infinite because of quantifier terms.

\begin{warning}[Calculus dependence]
Adding equality rules, induction, fixed points, modal rules, or infinitary rules
changes the cut reductions.  A theorem for the displayed \(\LK\) or \(\LJ\)
does not automatically apply to the extension.
\end{warning}

\begin{exercise}
\label{ex:cut-elimination-01}
Write the principal cut reduction for a cut formula \(A\land B\) when the left
premise ends in \(\land R\) and the right premise ends in \(\land L\).  Show
that each new cut has smaller cut rank.
\end{exercise}

\begin{exercise}
\label{ex:cut-elimination-02}
Explain why the subformula property gives a decision procedure for
propositional \(\LJ\), but not directly for first-order \(\LJ\).
\end{exercise}

